\documentclass[tikz,border=12pt]{standalone}
\usepackage{ctex}
\usepackage{tikz}
\usetikzlibrary{positioning,fit,shadows,backgrounds,calc}

\begin{document}
\begin{tikzpicture}[
  font=\small,
  % ---- 主流程模块：统一宽度与高度，等高由 minimum height 保证 ----
  stage/.style={
    draw=black, line width=0.5pt, rounded corners=3pt,
    fill=gray!18, minimum width=2.0cm, minimum height=2.9cm,
    align=center, font=\footnotesize, inner sep=4pt, text width=1.7cm
  },
  stagecore/.style={
    stage, fill=gray!25
  },
  % ---- 底部支撑条 ----
  bar/.style={
    draw=black, line width=0.5pt, rounded corners=2pt, fill=gray!10,
    minimum width=16.0cm, minimum height=0.6cm, align=center,
    font=\footnotesize, inner sep=3pt, text width=15.6cm
  },
  arrow/.style={->, >=stealth, line width=0.6pt},
  dasharrow/.style={->, >=stealth, line width=0.6pt, dashed},
  label/.style={font=\footnotesize}
]

% ==================== 主流程模块（链式 right=of 自动排布） ====================
\node[stage] (D) {
  \textbf{数据源}\\[2pt]
  案情库\\
  交易流水\\
  AMLSim\\
  合成数据
};
\node[stage, right=10pt of D] (P) {
  \textbf{预处理}\\[2pt]
  字段对齐\\
  数据脱敏\\
  格式转换\\
  案情结构化
};
\node[stagecore, right=10pt of P] (PL) {
  \textbf{规划}\\[2pt]
  （Plan）
};
\node[stage, right=10pt of PL] (A) {
  \textbf{分析}\\[2pt]
  （Analyze）\\[2pt]
  结构通道\\
  文本通道
};
\node[stagecore, right=10pt of A] (CL) {
  \textbf{聚类}\\[2pt]
  （Cluster）
};
\node[stagecore, right=10pt of CL] (R) {
  \textbf{反思}\\[2pt]
  （Reflect）
};
\node[stage, right=10pt of R] (APP) {
  \textbf{应用服务}\\[2pt]
  API 网关\\
  结果可视化\\
  人工审核\\
  审计日志
};

% ==================== 主数据流 ====================
\draw[arrow] (D.east) -- (P.west);
\draw[arrow] (P.east) -- (PL.west);
\draw[arrow] (PL.east) -- (A.west);
\draw[arrow] (A.east) -- (CL.west);
\draw[arrow] (CL.east) -- (R.west);
\draw[arrow] (R.east) -- (APP.west);

% ==================== 反思回环（下方留白区，基于节点相对坐标） ====================
\coordinate (rbotR) at ($(R.south)+(0,-1.25)$);
\coordinate (rbotL) at ($(PL.south)+(0,-1.25)$);
\draw[dasharrow] (R.south) -- (rbotR);
\draw[dasharrow] (rbotR) -- (rbotL)
  node[midway, below, yshift=-2pt, label] {置信度低 / 规则触发};
\draw[dasharrow] (rbotL) -- (PL.south)
  node[pos=0.5, left, xshift=-5pt, label] {返回规划};

% ==================== 人工反馈闭环（上方留白区，基于节点相对坐标） ====================
\coordinate (fbR) at ($(APP.north)+(0,1.25)$);
\coordinate (fbL) at ($(D.north)+(0,1.25)$);
\draw[dasharrow] (APP.north) -- (fbR);
\draw[dasharrow] (fbR) -- (fbL)
  node[midway, above, yshift=2pt, label] {人工审核反馈（Human-in-the-loop）};
\draw[dasharrow] (fbL) -- (D.north);

% ==================== 底部支撑保障（锚定整体水平中心） ====================
\coordinate (mid) at ($(D.west)!0.5!(APP.east)$);
\node[bar] (infra) at ([yshift=-3.0cm]mid)
  {数据与模型支撑：关系型数据库 · 缓存中间件 · 本地嵌入模型 · 云端大模型};
\node[bar] (trust) at ([yshift=-3.65cm]mid)
  {工程可信保障：数据不出域 · RBAC 鉴权 · 审计溯源 · 可解释性 · 多级降级};

\end{tikzpicture}
\end{document}
